\documentclass[12pt,a4paper]{article}

\usepackage[italian]{babel}
\usepackage[utf8]{inputenc}
\usepackage[T1]{fontenc}
\usepackage[expansion=false]{microtype}
\usepackage{geometry}
\geometry{
  top=2.5cm,
  bottom=2.5cm,
  left=2.5cm,
  right=2.5cm
}

\usepackage{setspace}
\setstretch{1.2}

\usepackage{hyperref}
\hypersetup{
  colorlinks=true,
  linkcolor=blue,
  urlcolor=blue
}

\setlength{\parindent}{0pt}
\setlength{\parskip}{6pt}

\begin{document}

\begin{center}
  {\Large\bfseries \{titoloDocumento\}}\\[0.5cm]
  {\large \{sottotitolo\}}\\[1cm]
\end{center}

\noindent
\textbf{Committente}\\
\{intestazioneCommittente\}\\[0.25cm]

\noindent
\textbf{Impianto}\\
\{dettaglioImpianto\}\\[0.25cm]

\noindent
\textbf{Data documento}\\
\{dataDocumento\}\\[1cm]

\section*{Nota}

Questo è un \textbf{template di prova} completamente indipendente
rispetto alla relazione tecnica principale.

Puoi modificare liberamente questa struttura LaTeX,
aggiungere sezioni, immagini e tabelle senza toccare il codice
del backend o del frontend, a patto di mantenere i placeholder
che ti servono (come \verb|{titoloDocumento}|, \verb|{intestazioneCommittente}|, ecc.).

\vfill
\begin{center}
  {\small Generato automaticamente dal generatore di template per documenti tecnici.}
\end{center}

\end{document}

